\section{Background}
The seismic performance of reinforced-concrete buildings in Nepal is influenced by tectonic hazard, local site conditions, structural configuration, detailing, and construction quality. Gautam et al. documented recurring structural and construction deficiencies in Nepalese buildings, including weaknesses that can promote undesirable seismic mechanisms \cite{gautam2016}. For a hospital, the consequence of such deficiencies is not limited to collapse prevention; excessive drift, foundation movement, and damage to structural or nonstructural components can also interrupt post-earthquake operation.

The prototype considered in this study is a G+5 reinforced-concrete hospital building modelled as a Special Moment Resisting Frame (SMRF). Seismic loading is defined in accordance with NBC 105:2025. A conventional fixed-base model is compared with flexible-base models representing hard, medium, and soft soil. Translational, rocking, and torsional foundation stiffnesses are used to represent the soil--foundation interface, following the rigid-foundation impedance framework described by Gazetas \cite{gazetas1991}.

The assessment extends beyond elastic response quantities. SAP2000 pushover capacity curves are examined in both principal horizontal directions. The reported nonlinear-static performance coordinates are interpreted using the FEMA~356 target-displacement framework so that the fixed and flexible support conditions can be compared on a common displacement-based basis. Preliminary fragility relationships are then derived from the nonlinear time-history results of twelve ground motions using the adopted drift-based damage-state thresholds.
